%! Solving Quadratics by Factoring — Unit 3, Lesson 3.2 — Guided Notes (Answer Key)
\documentclass[10pt]{article}
\usepackage{algebra2-article}
\usepackage{algebra2-key}   % pulls in -boxes; do NOT also load -boxes
\newcommand{\TallMath}[1]{$\displaystyle #1\rule[-1.4em]{0pt}{3.2em}$}
\newcommand{\lgrid}[4]{%
  \draw[step=1, linegray!40, very thin] (#1,#3) grid (#2,#4);
  \draw[->, semithick, charcoal] (#1,0)--(#2,0) node[below right, font=\scriptsize]{$x$};
  \draw[->, semithick, charcoal] (0,#3)--(0,#4) node[left, font=\scriptsize]{$y$};}
% Vocab answer for the key: bold forest term, then the filled definition on a write-line.
\newcommand{\vocabans}[2]{%
  \par\noindent\textbf{\textcolor{forest}{#1:}}\\[1pt]\ansline{#2}\par}

\begin{document}

\pageheader{Unit 3, Lesson 3.2}{Guided Notes --- Answer Key}

\vspace{0.10in}

\begin{objectivebox}
\textbf{By the end of this lesson, I will be able to\dots}
\begin{itemize}\setlength{\itemsep}{2pt}
  \item \emph{factor} a quadratic completely: GCF first, then $x^2+bx+c$ by product/sum, $ax^2+bx+c$ by grouping, and the special patterns;
  \item \emph{solve} $ax^2+bx+c=0$ with the \emph{Zero Product Property} and verify the roots;
  \item connect the \emph{roots} to the parabola's \emph{$x$-intercepts / zeros}.
\end{itemize}
\end{objectivebox}

\vspace{0.06in}

\begin{vocabbox}
\small
Fill in each term as we build it together.
\par\vspace{2pt}
\vocabans{Factoring}{writing a polynomial as a product of factors --- the reverse of multiplying (un-FOIL)}
\vocabans{Zero Product Property}{if $AB=0$ then $A=0$ or $B=0$; a product is zero only when a factor is zero}
\vocabans{Root / zero / solution}{three names for an $x$-value that makes $f(x)=0$ --- the parabola's $x$-intercepts}
\vocabans{Difference of squares}{$x^2-k^2=(x-k)(x+k)$}
\vocabans{Perfect-square trinomial (double root)}{$x^2\pm2kx+k^2=(x\pm k)^2$; factoring gives one repeated root where the graph touches the axis}
\end{vocabbox}

\begin{hookbox}
In Lesson~3.1 you were \textbf{handed} the factored form of the unit's anchor parabola:
\[ f(x)=x^2-2x-3 \;=\; (x+1)(x-3). \]
\begin{itemize}\setlength{\itemsep}{1pt}
  \item Where do you think the ``$+1$'' and ``$-3$'' came from? \ansline{Two integers with product $-3$ and sum $-2$: namely $+1$ and $-3$.}
  \item If $(x+1)(x-3)=0$, what must be true about \emph{one} of the factors? \ansline{One factor must equal $0$: $x+1=0$ or $x-3=0$.}
\end{itemize}
Today we \emph{produce} the factored form ourselves --- and use it to find where the graph meets the $x$-axis.
\end{hookbox}

\boxguard[20]
\begin{notesbox}{1. The Zero Product Property --- the engine}
If a product of two numbers is \textbf{zero}, then at least one of them must \emph{be} zero:
\[ \text{if } A\cdot B = 0, \text{ then } A=\textcolor{keyred}{\mathbf{0}} \ \text{ or } \ B=\textcolor{keyred}{\mathbf{0}}. \]
That is the \emph{only} way a product hits zero. So a factored equation becomes two easy linear ones.
\begin{center}
\begin{minipage}{0.54\linewidth}
\textbf{Solve} $(x+1)(x-3)=0$:
\begin{enumerate}[leftmargin=1.6em, itemsep=2pt]
  \item Set each factor to $0$: \quad $x+1=0$ \ or \ $x-3=0$.
  \item Solve each: \quad $x=\textcolor{keyred}{\mathbf{-1}}$ \ or \ $x=\textcolor{keyred}{\mathbf{3}}$.
\end{enumerate}
These are the \textbf{roots} --- and the parabola's \ans{$x$-intercepts}.

\vspace{3pt}
\textbf{Careful:} the property needs a \textbf{zero}. If $(x+1)(x-3)=\underline{\ 5\ }$, you may
\emph{not} write $x+1=5$. First move everything to one side.
\end{minipage}%
\begin{minipage}{0.44\linewidth}\centering
\begin{tikzpicture}[scale=0.42]
  \lgrid{-3}{5}{-5}{3}
  \begin{scope}\clip (-3,-5) rectangle (5,3);
    \draw[very thick, forest, domain=-1.4:3.4, samples=75, smooth] plot (\x,{(\x+1)*(\x-3)});
  \end{scope}
  \fill[charcoal] (-1,0) circle (3pt) node[above left, font=\scriptsize]{$-1$};
  \fill[charcoal] (3,0) circle (3pt) node[above right, font=\scriptsize]{$3$};
  \fill[charcoal] (1,-4) circle (3pt);
\end{tikzpicture}
\end{minipage}
\end{center}
\end{notesbox}

\begin{notesbox}{2. Factor $x^2+bx+c$ --- un-FOIL}
Since $(x+p)(x+q)=x^2+(p+q)x+pq$, factoring reverses the search: find two integers whose
\textbf{product is $c$} and whose \textbf{sum is $b$}.
\begin{center}
\renewcommand{\arraystretch}{1.4}
\begin{tabular}{|c|c|}
\hline
\textbf{sign of $c$} & \textbf{what the two integers look like} \\ \hline
$c>0$ & same sign (both match the sign of $b$) \\ \hline
$c<0$ & opposite signs \\ \hline
\end{tabular}
\end{center}
\textbf{Example (the anchor).} Factor $x^2-2x-3$: need product $\textcolor{keyred}{\mathbf{-3}}$, sum
$\textcolor{keyred}{\mathbf{-2}}$. The integers are $\textcolor{keyred}{\mathbf{+1}}$ and
$\textcolor{keyred}{\mathbf{-3}}$, so
\[ x^2-2x-3 = (\,x\ \textcolor{keyred}{\mathbf{+1}}\,)(\,x\ \textcolor{keyred}{\mathbf{-3}}\,). \]
\textbf{You try.} Factor $x^2+7x+12$: product $12$, sum $7 \Rightarrow \textcolor{keyred}{\mathbf{3}},\
\textcolor{keyred}{\mathbf{4}}$; so $x^2+7x+12 = \textcolor{keyred}{\mathbf{(x+3)(x+4)}}$.
\end{notesbox}

\begin{notesbox}{3. GCF first --- and the \#1 trap}
\emph{Always} pull out the greatest common factor before anything else.
\[ 2x^2-8x = \textcolor{keyred}{\mathbf{2x}}\,(\,x-\textcolor{keyred}{\mathbf{4}}\,). \]
\textbf{Solve} $2x^2-8x=0$: factor to $2x(x-4)=0$, so $x=\textcolor{keyred}{\mathbf{0}}$ or
$x=\textcolor{keyred}{\mathbf{4}}$.

\vspace{3pt}
\fbox{\parbox{0.96\linewidth}{\small\textbf{Trap:} to solve $3x^2=12x$, do \emph{not} divide both sides
by $x$ --- that deletes the root $x=0$. Instead move to one side: $3x^2-12x=0 \Rightarrow 3x(x-4)=0
\Rightarrow x=\textcolor{keyred}{\mathbf{0}}$ or $x=\textcolor{keyred}{\mathbf{4}}$.}}
\end{notesbox}

\begin{notesbox}{4. Two special patterns + $a\ne1$ by grouping}
\textbf{Difference of squares:} $\;x^2-k^2=(x-k)(x+k)$. \quad
Factor $x^2-16 = (\,x-\textcolor{keyred}{\mathbf{4}}\,)(\,x+\textcolor{keyred}{\mathbf{4}}\,)$; roots
$x=\textcolor{keyred}{\mathbf{\pm4}}$.

\vspace{2pt}
\textbf{Perfect-square trinomial:} $\;x^2-6x+9=(x-\textcolor{keyred}{\mathbf{3}})^2$; this gives one
\textbf{double root} $x=\textcolor{keyred}{\mathbf{3}}$ (the parabola \emph{touches} the $x$-axis there).

\vspace{2pt}
\textbf{Leading coefficient $a\ne1$ --- grouping ($a\!\cdot\!c$ method).} Factor $2x^2+7x+3$:
\begin{enumerate}[leftmargin=1.7em, itemsep=1pt]
  \item $a\!\cdot\!c = 2\cdot3 = \textcolor{keyred}{\mathbf{6}}$; two integers with that product and sum $7$: $\textcolor{keyred}{\mathbf{6}},\ \textcolor{keyred}{\mathbf{1}}$.
  \item Split the middle term: $2x^2+\textcolor{keyred}{\mathbf{6x}}+\textcolor{keyred}{\mathbf{x}}+3$.
  \item Group: $2x(x+3)+1(x+3) = (\,\textcolor{keyred}{\mathbf{2x+1}}\,)(\,x+3\,)$; roots $x=\textcolor{keyred}{\mathbf{-\tfrac12}},\ \textcolor{keyred}{\mathbf{-3}}$.
\end{enumerate}
\end{notesbox}

\boxguard[22]
\begin{practicebox}
\textbf{Solve each by factoring. Show the factored form, then the roots.}
\begin{enumerate}[leftmargin=1.9em, itemsep=6pt]
  \item $x^2-x-6=0$
    \begin{work}
      x^2-x-6 &= 0 \\
      (x-3)(x+2) &= 0 \\
      x-3 = 0 \quad&\text{or}\quad x+2 = 0 \\
      x &= 3,\ -2
    \end{work}
  \item $x^2-5x=0$
    \begin{work}
      x^2-5x &= 0 \\
      x(x-5) &= 0 \\
      x = 0 \quad&\text{or}\quad x-5 = 0 \\
      x &= 0,\ 5
    \end{work}
  \item $x^2-49=0$
    \begin{work}
      x^2-49 &= 0 \\
      (x-7)(x+7) &= 0 \\
      x-7 = 0 \quad&\text{or}\quad x+7 = 0 \\
      x &= 7,\ -7
    \end{work}
\end{enumerate}
\end{practicebox}

\end{document}
